% Generated by uscribe -- do not edit by hand
\documentclass[letterpaper,12pt]{article}
\usepackage[T1]{fontenc}
\usepackage[utf8]{inputenc}
\usepackage{graphicx}
\usepackage{hyperref}
\usepackage{xcolor}
\usepackage{fancyvrb}
\IfFileExists{fvextra.sty}{%
  \usepackage{fvextra}%
  \fvset{breaklines=true,breakanywhere=true}%
}{}
\usepackage{booktabs}
\usepackage{longtable}
\usepackage{array}
\IfFileExists{multirow.sty}{\usepackage{multirow}}{}
\IfFileExists{caption.sty}{%
  \usepackage{caption}%
  \newcommand{\uscribecaptionstar}[1]{\caption*{##1}}%
}{%
  \newcommand{\uscribecaptionstar}[1]{{\centering\small ##1\par}}%
}
\hypersetup{colorlinks=true,linkcolor=blue,urlcolor=blue}
\begin{document}
\begin{titlepage}
\null\vfil\vskip 40pt
\begin{center}
{\LARGE\bfseries Network Namespaces in Unicon \par}
\vskip 3em
{\large\bfseries Jafar Al-Gharaibeh \par}
\vskip 1.0em
{\bfseries Unicon Technical Report: 30\par}
\vskip 1.0em
{\large\bfseries July 2026\par}
\vskip 5.5em
{\bfseries Abstract}\par
\begin{quote}
Linux network namespaces isolate a process's network stack -- interfaces, addresses, and routes -- without a full container. This report describes native support in Unicon: open() mode "j" creates a namespace handle; an optional leading handle on fork(), spawn(), and system() runs the child inside that namespace; ns["name"] / key(ns) peek status. Trailing persist=, bridge=, and userns= attributes cover ip(8) visibility, a pre-created bridge, and unprivileged bootstrap. Package net supplies link helpers (veth, bridge) on top of those handles.
\end{quote}
\begin{quote}
{\bfseries Keywords:} Unicon, network namespaces, Linux, open, fork, spawn, system, runtime, technical report.
\end{quote}
\vskip 0.8in
{\large\bfseries
Unicon Project\\
\url{http://unicon.org}\\
\ \
University of Idaho\\
Department of Computer Science\\
Moscow, ID, 83844, USA
}
\end{center}
\vfil\null
\end{titlepage}
\setcounter{page}{1}
\section{1. Introduction}
\label{utr30-sec-intro}

This report describes Linux network namespaces [15] in Unicon: how a handle is opened, how a child process or thread joins it, how status is peeked, and what remains out of scope. It is formatted as a Unicon Technical Report [1]. The language-facing reference lives in \emph{Programming with Unicon} [7] (Chapter 5, "Network Namespaces"; Chapter 14, \texttt{unettopo}) and the language reference (\texttt{open} mode \texttt{j}, optional leading namespace file on \texttt{fork} / \texttt{spawn} / \texttt{system}, subscript peek, \texttt{key()}). Automated tests live in \texttt{tests/\allowbreak{}posix/\allowbreak{}netns.\allowbreak{}icn} and \texttt{tests/\allowbreak{}unicon/\allowbreak{}nettopo.\allowbreak{}icn}.


The feature is optional and Linux-only. A build with \texttt{unshare()} / \texttt{setns()} and \texttt{CLONE\_\allowbreak{}NEWNET} reports \texttt{Network~namespaces} in \texttt{\&features} (the \texttt{\_\allowbreak{}NETNS} flag). Other hosts omit the feature; mode \texttt{j} fails with error 1045.

\section{2. Motivation}
\label{utr30-sec-motivation}

Test topologies on Linux are often a handful of network namespaces, veth pairs, and optional bridges: isolated stacks wired together, cheap enough to create per test run. Unicon already treats \texttt{open()} as the constructor for files, sockets, TLS, SSH, and crypto handles. Namespace membership is the same kind of resource: create it, pass it, close it. Shelling out to \texttt{ip~netns} works, but it cannot join the Unicon process itself, and it cannot put \texttt{fork()} / \texttt{spawn()} / \texttt{system()} children into a namespace without a second \texttt{ip~netns~exec} wrapper.

\section{3. Design principles}
\label{utr30-sec-principles}

\textbf{One OS thread is in one network namespace for its lifetime.} Membership lives on the kernel \texttt{nsproxy}. \texttt{fork(ns)} and \texttt{system(ns,\allowbreak{}~.\allowbreak{}.\allowbreak{}.\allowbreak{})} join in the child process; \texttt{spawn(ns,\allowbreak{}~.\allowbreak{}.\allowbreak{}.\allowbreak{})} joins on the new OS thread before any Icon code runs. The runtime never migrates a live thread between namespaces.


\textbf{Existing verbs, not new globals.} Mode \texttt{j} and trailing \texttt{name=value} attributes follow SSL and SSH. \texttt{fork}, \texttt{spawn}, and \texttt{system} take an optional leading namespace file, the same role a window plays in \texttt{DrawPoint(w,\allowbreak{}~x,\allowbreak{}~y)}. Link configuration that \texttt{ip(8)} already expresses stays in \texttt{package~net}.


\textbf{Status is peeked; \texttt{Attrib()} does not set netns fields.} \texttt{ns["name"]} and \texttt{key(ns)} match sockets, TLS, SSH, and crypto handles. Writes for interfaces, addresses, and routes are deferred; use \texttt{package~net}.


\textbf{Create needs privilege or a user namespace.} \texttt{unshare(CLONE\_\allowbreak{}NEWNET)} requires \texttt{CAP\_\allowbreak{}NET\_\allowbreak{}ADMIN}, or a successful unprivileged user-namespace bootstrap (\hyperref[utr30-sec-userns]{section 6}).

\section{4. Opening and closing}
\label{utr30-sec-open}

Mode \texttt{"j"} creates a namespace named \texttt{s1}. The result is a file handle, not a read/write stream. Other mode letters combined with \texttt{j} are error 209.

\begin{Verbatim}[fontsize=\small,frame=single,commandchars=\\\{\}]
procedure main()
   local ns
   ns := open("r1", "j") | stop(&errortext)
   write(image(ns))
   close(ns)
end
\end{Verbatim}


Output:

\begin{Verbatim}[fontsize=\small,frame=single,commandchars=\\\{\}]
netns(r1)
\end{Verbatim}


Implementation: a short-lived child calls \texttt{unshare(CLONE\_\allowbreak{}NEWNET)}, bind-mounts \texttt{/\allowbreak{}proc/\allowbreak{}self/\allowbreak{}ns/\allowbreak{}net} onto a private path under \texttt{/\allowbreak{}tmp/\allowbreak{}unicon-\allowbreak{}netns-\allowbreak{}<uid>/\allowbreak{}}, brings \texttt{lo} up, and exits. The private mount pins the kernel namespace. \texttt{Fs\_\allowbreak{}NETNS} (\texttt{040} when the feature is built) marks the handle; state is a \texttt{struct~NetnsFile~*} on the file block.


Trailing attributes are \texttt{name=value} strings. Booleans are exactly \texttt{yes} or \texttt{no}. Unknown names or other boolean spellings are error 205.

\begin{longtable}{>{\raggedright\sloppy\arraybackslash}p{\dimexpr(\textwidth-2\tabcolsep*2)*1/2+2\tabcolsep*(1-1)\relax}>{\raggedright\sloppy\arraybackslash}p{\dimexpr(\textwidth-2\tabcolsep*2)*1/2+2\tabcolsep*(1-1)\relax}}
\hline
Attribute & Role \\
\hline
\endfirsthead
\hline
Attribute & Role \\
\hline
\endhead
\hline
\endfoot
\hline
\endlastfoot
\texttt{persist=yes|no} & Also bind-mount at \texttt{/\allowbreak{}var/\allowbreak{}run/\allowbreak{}netns/\allowbreak{}<name>} so \\
\texttt{bridge=yes|no} & Best-effort \texttt{bridge0} inside the new namespace. \\
\texttt{userns=yes|no} & Force or forbid the unprivileged bootstrap in \\
\end{longtable}


Ephemeral opens (no \texttt{persist}) have no user-visible \texttt{ip~netns} name. Two ephemeral \texttt{open("r1",\allowbreak{}~"j")} calls yield two independent objects. \texttt{persist=yes} uses \texttt{O\_\allowbreak{}EXCL}; a second persist open of the same name fails.


\texttt{close(ns)} always unmounts this handle's pins. It does not wait on \texttt{ns["refcount"]}. Unmounting a pin does not destroy the kernel namespace: any process that already joined still holds an \texttt{nsproxy} reference. After \texttt{close()}, a persist name disappears from \texttt{ip~netns~list}, and a later \texttt{open()} of that name creates a new namespace. Live handles are also unmounted at process exit and on \texttt{SIGINT} / \texttt{SIGTERM} when those signals still have the default disposition. Cleanup runs only in the creating process, so \texttt{fork()} children that exit do not tear down the parent's pins. \texttt{SIGKILL} can still leave orphan mounts.

\section{5. Joining a namespace}
\label{utr30-sec-exec}

The namespace file leads when it is the scope of the call.


\textbf{\texttt{fork(ns)}} -- POSIX \texttt{fork()} with one optional argument. In the child, \texttt{setns()} runs before the function returns, so the first Icon statement is already namespaced. \texttt{fork()} with no argument is unchanged. \texttt{exec()} is unchanged: a child that already joined does not need extra namespace awareness.


\textbf{\texttt{spawn(ns,\allowbreak{}~ce,\allowbreak{}~.\allowbreak{}.\allowbreak{}.\allowbreak{})}} -- optional leading file, then the existing block/string/stack sizes. \texttt{spawn()} holds the namespace object, stores it on the co-expression, and starts the OS thread. The trampoline joins before any Icon code runs on that thread, then releases the hold, so \texttt{close(ns)} from the parent cannot race the join. With concurrent threads, the OS thread starts immediately; no \texttt{@\allowbreak{}ce} activation is required. \texttt{wait(ce)} waits for it to finish.


\textbf{\texttt{system(ns,\allowbreak{}~argv,\allowbreak{}~.\allowbreak{}.\allowbreak{}.\allowbreak{})}} -- same leading-file convention. The child joins before \texttt{dup} / \texttt{exec}.

\begin{Verbatim}[fontsize=\small,frame=single,commandchars=\\\{\}]
procedure main()
   local ns, pid
   ns := open("r1", "j") | stop(&errortext)
   pid := fork(ns) | stop("fork failed: ", &errortext)
   if pid = 0 then \{
      system(["ip", "link", "show"])
      exit(0)
      \}
   wait(pid)
   system(ns, ["ip", "link", "show"])
   close(ns)
end
\end{Verbatim}


\texttt{fork()} fails (no value) only on the system-call error. A successful child result is \texttt{0}, so discriminate with comparison, not \texttt{else}.

\section{6. Unprivileged bootstrap}
\label{utr30-sec-userns}

Creating a network namespace needs \texttt{CAP\_\allowbreak{}NET\_\allowbreak{}ADMIN} in the current user namespace [16]. \texttt{userns=yes} enters a new user namespace, maps the calling uid/gid to 0 inside it, enters a private mount namespace, and then creates the netns. The mapping is written by a short-lived \texttt{fork()} helper that stays in the host user namespace; writing \texttt{uid\_\allowbreak{}map} from the same process that called \texttt{unshare(CLONE\_\allowbreak{}NEWUSER)} is rejected on hosts that restrict unprivileged user namespaces.


If \texttt{userns=} is omitted and the privileged create fails with \texttt{EPERM} or \texttt{EACCES}, \texttt{open(.\allowbreak{}.\allowbreak{}.\allowbreak{},\allowbreak{}~"j")} retries once after the same bootstrap. Explicit \texttt{userns=yes} bootstraps up front. Explicit \texttt{userns=no} never bootstraps.


\texttt{unshare(CLONE\_\allowbreak{}NEWUSER)} fails with \texttt{EINVAL} if the process is already multithreaded. That must happen before any \texttt{spawn()}-created thread exists; the \texttt{\&errortext} states the constraint.


After a successful bootstrap the Unicon process is uid 0 \emph{inside} the user namespace for the rest of its life. Later \texttt{open(.\allowbreak{}.\allowbreak{}.\allowbreak{},\allowbreak{}~"j")}, \texttt{fork(ns)}, \texttt{spawn(ns,\allowbreak{}~.\allowbreak{}.\allowbreak{}.\allowbreak{})}, \texttt{system(ns,\allowbreak{}~.\allowbreak{}.\allowbreak{}.\allowbreak{})}, and \texttt{ip} via \texttt{package~net} need no further host privilege. A private tmpfs is mounted on \texttt{/\allowbreak{}var/\allowbreak{}run/\allowbreak{}netns} so \texttt{persist=yes} can create entries when the host directory is not writable to a mapped-root user namespace.


Linux security modules may still deny further \texttt{CLONE\_\allowbreak{}NEWNS} / \texttt{CLONE\_\allowbreak{}NEWNET} unshares for unprivileged callers. In that case \texttt{open()} fails with the system \texttt{\&errortext}; run with \texttt{CAP\_\allowbreak{}NET\_\allowbreak{}ADMIN}, or relax the host policy.


\texttt{ns["userns"]} reports whether this handle was created after bootstrap. The choice is fixed at \texttt{open()}.

\section{7. Status peek}
\label{utr30-sec-peek}
\begin{Verbatim}[fontsize=\small,frame=single,commandchars=\\\{\}]
procedure main()
   local ns, k
   ns := open("r1", "j", "persist=yes") | stop(&errortext)
   write("name=", ns["name"])
   write("persist=", \\ns["persist"] | "no")
   write("userns=", \\ns["userns"] | "no")
   write("refcount=", ns["refcount"])
   every k := key(ns) do
      write("field ", k)
   close(ns)
end
\end{Verbatim}


Unknown names raise 1336 (\texttt{bad~netns~status~field}). An unpopulated field fails. Boolean fields that answered succeed with \texttt{"yes"} or \texttt{\&null}. \texttt{key(ns)} generates every answerable field. \texttt{ns["*"]} snapshots those fields under one lock. Peeking a closed handle is error 174. \texttt{key(ns)} that has not yet produced a name also raises 174 if the handle is already closed.

\begin{longtable}{>{\raggedright\sloppy\arraybackslash}p{\dimexpr(\textwidth-2\tabcolsep*2)*1/2+2\tabcolsep*(1-1)\relax}>{\raggedright\sloppy\arraybackslash}p{\dimexpr(\textwidth-2\tabcolsep*2)*1/2+2\tabcolsep*(1-1)\relax}}
\hline
Name & Value \\
\hline
\endfirsthead
\hline
Name & Value \\
\hline
\endhead
\hline
\endfoot
\hline
\endlastfoot
\texttt{name} & Namespace name from \texttt{open()} \\
\texttt{persist} & \texttt{"yes"} or \texttt{\&null} \\
\texttt{userns} & \texttt{"yes"} or \texttt{\&null} \\
\texttt{refcount} & Integer; the open file plus any \texttt{spawn()} hold \\
\end{longtable}


\texttt{Attrib()} has no netns setters. Interface, address, and route configuration lives in \texttt{package~net}.

\section{8. Link helpers}
\label{utr30-sec-pkg}

\texttt{uni/\allowbreak{}lib/\allowbreak{}netns.\allowbreak{}icn} is pure Unicon. It shells out to \texttt{ip} / \texttt{ip~netns}. There is no rtnetlink binding.


\texttt{Netns(name)} prefers native \texttt{open(name,\allowbreak{}~"j",\allowbreak{}~"persist=yes")} and keeps the file open so the persist mount stays pinned. If that fails it tries an ephemeral native open, then \texttt{ip~netns~add}. \texttt{delete()} closes the file or runs \texttt{ip~netns~del}.


\texttt{Link} is the interface base. \texttt{Veth}, \texttt{Bridge}, \texttt{Vrf}, and \texttt{Tun} create devices; a bare \texttt{Link(name)} wraps an existing interface. Every mutating method goes through \texttt{run()}, which prefers \texttt{system(nsfile,\allowbreak{}~argv)} on a native handle and otherwise uses \texttt{ip~netns~exec}.


\texttt{Link.\allowbreak{}move(target)} accepts a \texttt{Netns} instance, a string name, or a native file (\texttt{target["name"]}). \texttt{Veth.\allowbreak{}peer()} returns a host-resident \texttt{Link} for the other end and does not copy \texttt{ns}; the two ends usually move into different namespaces. Create a veth \emph{inside} a namespace when running unprivileged: host RTNETLINK create is often denied after user-namespace bootstrap.

\begin{Verbatim}[fontsize=\small,frame=single,commandchars=\\\{\}]
import net

procedure main()
   local a, b, v, peer
   a := Netns("r1") | stop(&errortext)
   b := Netns("r2") | stop(&errortext)
   v := Veth("v1", "v2", a) | stop(&errortext)
   peer := v.peer()
   peer.move(b) | stop(&errortext)
   v.addr("10.0.0.1/24")
   peer.addr("10.0.0.2/24")
   v.up()
   peer.up()
   system(a.nsfile, ["ping", "-c", "1", "10.0.0.2"])
   v.delete()
   a.delete()
   b.delete()
end
\end{Verbatim}


\texttt{uni/\allowbreak{}progs/\allowbreak{}unettopo.\allowbreak{}icn} is that topology as a demo.

\section{9. Runtime implementation}
\label{utr30-sec-impl}

\textbf{Open / close.} \texttt{src/\allowbreak{}runtime/\allowbreak{}fsys.\allowbreak{}r} parses mode \texttt{j} and the trailing attributes, then calls \texttt{netns\_\allowbreak{}create()}. \texttt{close()} releases the \texttt{NetnsFile}.


\textbf{Join.} \texttt{src/\allowbreak{}runtime/\allowbreak{}fxposix.\allowbreak{}ri} (\texttt{fork}), \texttt{src/\allowbreak{}runtime/\allowbreak{}fsys.\allowbreak{}r} (\texttt{system}), and \texttt{src/\allowbreak{}runtime/\allowbreak{}rcoexpr.\allowbreak{}r} (\texttt{nctramp} for \texttt{spawn}) call \texttt{netns\_\allowbreak{}join()}. \texttt{spawn()} in \texttt{fmisc.\allowbreak{}r} uses \texttt{OptNetns} (\texttt{grttin.\allowbreak{}h}), the same optional-leading-arg pattern as \texttt{OptWindow}.


\textbf{Helpers.} \texttt{src/\allowbreak{}common/\allowbreak{}rnetns.\allowbreak{}c} implements create, join, hold/release, userns bootstrap, and the exit/signal reaper. \texttt{configure.\allowbreak{}ac} link-checks \texttt{unshare} / \texttt{setns} / \texttt{CLONE\_\allowbreak{}NEWNET}. \texttt{Fs\_\allowbreak{}NETNS} is \texttt{0} when the feature is off so \texttt{OptNetns} still compiles.


\textbf{Peek.} \texttt{netns\_\allowbreak{}peek} in \texttt{src/\allowbreak{}runtime/\allowbreak{}rposix.\allowbreak{}r} uses the same \texttt{filepeek} table as sockets and SSH. \texttt{oref.\allowbreak{}r} implements \texttt{ns["name"]}; \texttt{fstruct.\allowbreak{}r} implements \texttt{key(ns)}. Unknown names are 1336.


\textbf{Image.} \texttt{image(ns)} is \texttt{netns(<name>)}.


\textbf{CLOEXEC.} No long-lived namespace fd is held; pins are paths. \texttt{netns\_\allowbreak{}join} opens with \texttt{O\_\allowbreak{}CLOEXEC}. Create and bootstrap pipes use \texttt{pipe2(O\_\allowbreak{}CLOEXEC)}.

\section{10. Status and remaining work}
\label{utr30-sec-status}

The language surface described here is implemented: mode \texttt{j}, persist/bridge/userns attributes, \texttt{fork(ns)} / \texttt{spawn(ns,\allowbreak{}~.\allowbreak{}.\allowbreak{}.\allowbreak{})} / \texttt{system(ns,\allowbreak{}~.\allowbreak{}.\allowbreak{}.\allowbreak{})}, status peek, \texttt{package~net} veth/bridge helpers, and the \texttt{unettopo} demo.


\textbf{Native rtnetlink.} Link helpers still shell out to \texttt{ip(8)}.


\textbf{Other platforms.} No BSD jail or Windows backend.


\textbf{PID and mount isolation.} \texttt{fork(ns)} joins only the network namespace.


\textbf{\texttt{Attrib()} writes} for interface, address, and route state. Use \texttt{package~net}.


\textbf{\texttt{Netns.\allowbreak{}list()} / \texttt{Netns.\allowbreak{}current()}}, link statistics, and macvlan/ipvlan classes are not provided.

\section{Appendix: test coverage}
\label{utr30-sec-coverage}

\texttt{tests/\allowbreak{}posix/\allowbreak{}netns.\allowbreak{}icn} drops to \texttt{nobody} when started as root, then opens with \texttt{userns=yes}. It checks:

\begin{itemize}
\item 

feature string and a successful \texttt{open(.\allowbreak{}.\allowbreak{}.\allowbreak{},\allowbreak{}~"j")}


\item 

\texttt{ns["name"]}, boolean \texttt{persist} / \texttt{userns}, \texttt{refcount}, \texttt{key(ns)}, and \texttt{ns["*"]}


\item 

\texttt{fork(ns)} and \texttt{system(ns,\allowbreak{}~.\allowbreak{}.\allowbreak{}.\allowbreak{})} see only \texttt{lo}


\item 

\texttt{spawn(ns,\allowbreak{}~.\allowbreak{}.\allowbreak{}.\allowbreak{})} when concurrent threads are present


\item 

persist listing in \texttt{ip~netns~list}, gone after \texttt{close()}


\item 

ephemeral name reuse vs persist \texttt{O\_\allowbreak{}EXCL}


\end{itemize}

Without a working userns or \texttt{CAP\_\allowbreak{}NET\_\allowbreak{}ADMIN} it prints \texttt{This~Program~Requires~.\allowbreak{}.\allowbreak{}.\allowbreak{}} so \texttt{Makefile.\allowbreak{}test} does not treat the run as \texttt{FAIL}.


\texttt{tests/\allowbreak{}unicon/\allowbreak{}nettopo.\allowbreak{}icn} covers \texttt{package~net}: two \texttt{Netns} objects, a veth created inside the first namespace, peer move, addresses, and \texttt{ip~netns~exec} iface lists.


Expected output is \texttt{tests/\allowbreak{}posix/\allowbreak{}stand/\allowbreak{}netns.\allowbreak{}std}:

\begin{Verbatim}[fontsize=\small,frame=single,commandchars=\\\{\}]
feature=yes
open=ok name=ut-ns persist=no userns=yes refcount=1
key=ok
fork-child: ifaces=lo
fork-parent: ok
system: ifaces=lo
spawn: ifaces=lo
close=ok
persist=ok name=ut-nsp persist=yes
bridge: ifaces=bridge0,lo
listed=yes
listed-after-close=no
dup-ephemeral=ok
dup-persist=eexist
done
\end{Verbatim}

\section{References}
\label{utr30-references}
\begin{thebibliography}{99}
\bibitem[1]{cite-jeffery-utr15}
Clinton Jeffery. "How to Write a Unicon Technical Report". doc/utr/utr15.tex. 2013.

\bibitem[2]{cite-libssh}
The libssh Project. "libssh -- The SSH Library". https://www.libssh.org/. 2026.

\bibitem[3]{cite-jeffery-utr13}
Clinton Jeffery. "The Unicon Messaging Facilities". doc/utr/utr13.odt.

\bibitem[4]{cite-openssl}
The OpenSSL Project. "OpenSSL -- Cryptography and SSL/TLS Toolkit". https://www.openssl.org/. 2026.

\bibitem[5]{cite-rfc2104}
H. Krawczyk and M. Bellare and R. Canetti. "HMAC: Keyed-Hashing for Message Authentication". RFC 2104. 1997.

\bibitem[6]{cite-rfc1112}
S. Deering. "Host Extensions for IP Multicasting". RFC 1112. 1989.

\bibitem[7]{cite-jeffery-pwu}
Clinton Jeffery and Shamim Mohamed and Jafar Al-Gharaibeh. "Programming with Unicon". doc/book/. 2026.

\bibitem[8]{cite-rfc4607}
H. Holbrook and B. Cain. "Source-Specific Multicast for IP". RFC 4607. 2006.

\bibitem[9]{cite-algharaibeh-utr29}
Jafar Al-Gharaibeh. "Raw Sockets and Packet Layouts in Unicon". doc/utr/utr29.rst. 2026.

\bibitem[10]{cite-algharaibeh-utr28}
Jafar Al-Gharaibeh. "Multicast and Socket Attributes in Unicon". doc/utr/utr28.rst. 2026.

\bibitem[11]{cite-rfc792}
J. Postel. "Internet Control Message Protocol". RFC 792. 1981.

\bibitem[12]{cite-rfc2236}
W. Fenner. "Internet Group Management Protocol, Version 2". RFC 2236. 1997.

\bibitem[13]{cite-rfc3376}
B. Cain and S. Deering and I. Kouvelas and B. Fenner and A. Thyagarajan. "Internet Group Management Protocol, Version 3". RFC 3376. 2002.

\bibitem[14]{cite-rfc1071}
R. Braden and D. Borman and C. Partridge. "Computing the Internet Checksum". RFC 1071. 1988.

\bibitem[15]{cite-namespaces7}
"namespaces(7)". Linux man-pages. 2024.

\bibitem[16]{cite-user-namespaces7}
"user\textbackslash{}\_namespaces(7)". Linux man-pages. 2024.

\end{thebibliography}
\end{document}
